\section{Methodology}
\label{sec:methodology}

The constrained learning framework can be instantiated in many ways depending on the function class $\mc{F}$ for outcome models. %
We 
present approximate solution methods to the constrained optimization problem~\eqref{eq:c-learner}
for linear models, gradient boosted regression trees, and deep neural networks. We empirically demonstrate these instantiations in Section~\ref{sec:experiments}.\footnote{Other frameworks~\citep{Nabi2024StatisticalLF} exist for constructing constrained outcome models, but do not restrict outcome models to $\mathcal F$; we exclude these from our analysis.}

\subsection{Data Splitting}
\label{sec:data_splitting}
We recommend sample splitting, where we split the data so
that nuisance estimators (e.g. $\what \pi, \what \mu, \what \mu^C$) are fitted on a training (auxiliary) fold and evaluated to form a
causal estimator on an evaluation (main) fold. 
Let $P_{\rm train}$ be the split on which nuisance parameters
such as the outcome model or the propensity score $\what{\pi}(\cdot)$ are trained, and let $P_{\rm val}$ be the split on which we perform model selection on them. 
We use a separate split $P_{\rm eval}$ to evaluate
our final causal estimators. %
Our constrained learning framework~\eqref{eq:c-learner} can thus be rewritten as
\vspace{-.2cm}
\begin{align}
& \what{\psi}^{\textrm {C-Learner}} := 
P_{\rm eval}\left[
\what{\mu}^C(X)
\right]~~~~~\mbox{where}  \nonumber\\
\label{eq:c-learner-datasets}
& \what{\mu}^C\in \argmin_{\wt\mu\in\mathcal F} \left\{
P_{\rm train}[ A (Y - \wt{\mu}(X))^2] : P_{\rm eval}
\left[\frac{A}{\what{\pi}(X)}
(Y-\wt\mu(X))\right]=0
\right\}.
\end{align}
For simplicity, we let $P_{\rm eval}=P_{\rm val}$ except where otherwise specified. 

\paragraph{Cross-fitting with $K$ folds}
\emph{Cross-fitting} with $K$ folds \citep{ChernozhukovChDeDuHaNeRo18} refers to the following sample splitting setup designed to utilize all of the data: 
first, we split the data into $K$ folds. Then, to evaluate on the $k$th fold in $P_{\rm eval}[\what \mu^C(X)]$, we train nuisance parameters on all but the $k$th fold in the data. We repeat this for all $K$ folds and average the results, so that the final estimator utilizes model evaluations over the entire dataset. 
Assuming $K$ evenly divides $n$ for simplicity, let $P_{k,n}$ be the empirical measure over data from the $k$-th fold, and $P_{-k,n}$ be the empirical measure over data from all other folds. 
Let $P_{\rm train}=P_{-k,n}$, and $P_{\rm val}=P_{\rm eval}=P_{k,n}$. 
For each $k=1,\ldots,K$, on the training fold $P_{-k, n}$, we train a propensity score model $\what{\pi}_{-k,n}(x)$ to estimate $P[A\mid X=x]$, as well as an unconstrained outcome model $\what{\mu}_{-k,n}(x)$, and a constrained outcome model $\what{\mu}_{-k,n}^C(x)$ to estimate $P[Y\mid A=1,X=x]$.  

Then, the AIPW and TMLE estimators from \Cref{sec:background} can be written as 
\begin{align*}
   \what\psi^{\rm AIPW}_{k,n}&=P_{k,n}[\what\mu_{-k,n}(X)]+P_{k,n}\left[\frac{A}{\pi_{-k,n}(X)}(Y-\what\mu_{-k,n}(X))\right]\\ 
   \what\psi^{\rm TMLE}_{k,n}&=P_{k,n}[\what\mu_{-k,n}(X)]+\frac{P_{k,n}[1/\pi_{-k,n}(X)]}{P_{k,n}[A/\pi_{-k,n}(X)^2]}P_{k,n}\left[\frac{A}{\pi_{-k,n}(X)}(Y-\what\mu_{-k,n}(X))\right].
\end{align*}

The C-Learner can be written as follows: the constrained outcome model optimizes the prediction loss evaluated under $P_{-k, n}$,  subject to the first-order correction constraint evaluated on the evaluation fold $P_{k, n}$, with final estimator $\what \psi_n^C$:
\vspace{-0.3em}
\begin{align}
\label{eq:constraint}
\what{\mu}_{-k, n}^C 
&\in \argmin_{\tilde{\mu} \in \mc{F}}
\left\{ P_{-k, n} [ A (Y - \tilde{\mu}(X))^2]:
P_{k,n}\left[\frac{A}{\what\pi_{-k,n}(X)}(Y-\tilde{\mu}(X))\right]=0
\right\}\\
    \what\psi^C_n &:=\frac{1}{K}\sum_{k=1}^K P_{k,n}[\what\mu_{-k,n}^C(X)].
\label{eq:c_learner_def}
\end{align}

In contrast to standard cross-fitting, in which nuisance parameters are learned only on the training split $P_{-k,n}$, 
the constraint~\eqref{eq:constraint} is over the eval split $P_{k,n}$. 

\subsection{Linear Models}
\label{sec:methods_linear}
When outcome models are linear functions of $X$, the
constrained learning problem~\eqref{eq:c-learner-datasets} 
has an analytic solution.
Using $\vec{\cdot}$ to denote stacked observations, define
\vspace{-0.3em}
\begin{equation*}
\begin{bmatrix}
\vec{Y}_{\rm train}  \\
\vec{X}_{\rm train} \\
\vec{H}_{\rm train} 
\end{bmatrix}
:= \begin{bmatrix}
\{Y_i\}_{i \in \mc{I}_{\rm train}}  \\
\{X_i\}_{i \in \mc{I}_{\rm train}} \\
\left\{ \frac{A_i}{\what{\pi}(X_i)} 
\right\}_{i \in \mc{I}_{\rm train}}
\end{bmatrix}
\qquad \mbox{and} \qquad
\begin{bmatrix}
\vec{Y}_{\rm eval}  \\
\vec{X}_{\rm eval} \\
\vec{H}_{\rm eval} 
\end{bmatrix}
:= \begin{bmatrix}
\{Y_i\}_{i \in \mc{I}_{\rm eval}}  \\
\{X_i\}_{i \in \mc{I}_{\rm eval}} \\
\left\{ \frac{A_i}{\what{\pi}(X_i)} 
\right\}_{i \in \mc{I}_{\rm eval}}
\end{bmatrix}
\end{equation*}
where $\mc{I}_{\rm train} = \{i \in \mbox{train}: A_i = 1\}$
and $\mc{I}_{\rm eval} = \{i \in \mbox{eval}: A_i = 1\}$ 
are indices with observations in each data split.
The constrained learning problem~\eqref{eq:c-learner-datasets} 
can be rewritten as \vspace{-1.5em}\begin{align*}
\what\theta^C 
    = \argmin_\theta \left\{
\frac{1}{2}\|\vec{Y}_{\rm train} - \vec{X}_{\rm train} \theta\|^2:\; 
\vec{H}_{\rm eval}^\top (\vec{Y}_{\rm eval} - \vec{X}_{\rm eval} \theta)=0 \right\}.
\end{align*}
The KKT conditions characterize the primal-dual optimum $(\what{\theta}^C, \what{\lambda})$
\vspace{-0.7em}
$$
\what{\theta}^C = (\vec{X}_{\rm train}^\top \vec{X}_{\rm train})^{-1} \vec{X}_{\rm train}^\top (\vec{Y}_{\rm train}+ \what{\lambda} \vec{H}_{\rm train})
~~\mbox{where}~~
\what{\lambda} = \frac{\vec{H}_{\rm eval}^\top
(\vec{Y}_{\rm eval} - \vec{Y}_{\rm ols})}{\vec{H}_{\rm eval}^\top \vec{X}_{\rm eval}
(\vec{X}_{\rm train}^\top \vec{X}_{\rm train})^{-1}
\vec{X}_{\rm train}^\top \vec{H}_{\rm train}}
$$
and $\vec{Y}_{\rm ols} := \vec{X}_{\rm train} (\vec{X}_{\rm train}^\top \vec{X}_{\rm train})^{-1}
\vec{X}_{\rm train}^\top \vec{Y}_{\rm train}$. %
Note that $\what \theta^C$ is the OLS with respect to the pseudo-label $\vec{Y}_{\rm train}+ \what{\lambda} \vec{H}_{\rm train}$ and the dual variable shifts the observed outcomes in the direction of $\vec{H}_{\rm train}$ similar to targeting~\eqref{eqn:targeting}, but with additional reweighting using covariates. 
When using our framework restricted to linear function classes, we obtain a new estimator that, to the best of our knowledge, cannot be recovered by existing methods. In \Cref{sec:ks}, we demonstrate that this new estimator improves upon existing methods.


\subsection{Gradient Boosted Regression Trees}
\label{sec:methods_boosting}
We consider outcome models that are gradient boosted regression trees. %
The gradient boosting framework~\cite{friedman2001greedy} iteratively estimates the functional gradient $g_j$ 
of the loss function evaluated on the current function estimate $\what\mu_{j}$.
For the standard  MSE loss $\ell(\mu;X,A,Y):=A(Y-\mu(X))^2$, 
we use weak learners in $\mc{G}$ (e.g., shallow decision trees) to compute
$$
\what{g}_{j+1} \in \argmin_{g \in \mathcal G} P_{\rm train} [(A(g_{j}(X,Y)-g(X)))^2]~~~\mbox{where}~~~g_{j}(X,Y;\what\mu_j) := \frac{\partial}{\partial \mu}\ell(\mu;X,Y)\Big|_{\mu = \what\mu_{j}}
$$
and set $\what \mu_{j+1} = \what\mu_j - \eta \what{g}_{j+1}$ for some step size $\eta$. 
For notational simplicity, since we assume $Y(0):=0$, let $\mu(X)$, $g_j(X,Y)$, $g(X,Y)$
be shorthand for $\mu(A,X)$, $g_j(X,A,Y)$, $g(X,A,Y)$ when $A=1$, and we let $\mu(0,X),g_j(0,X),g(0,X)=0$. 
The process repeats $J$ times until a maximum number of weak learners are fitted or an early stopping criterion is met. 

\begin{algorithm}[t]
\caption{C-Learner with Gradient Boosted Regression Trees}
\begin{algorithmic}[1]
\State \textbf{Input:} learning rate $\eta$, max trees $J$ and $K$, $\what \mu_0 := 0$
\For{$j = 0, 1, \dots, J-1$}
    \State Modify functional gradient $g_j = Y - \what{\mu}_j(X)$ to $\wt{g}_j = g_j + \epsilon^\star_j \cdot 1/\what{\pi}(X)$ where 
  $\epsilon_j^\star := \tfrac{P_{\rm train} [(Y - \what{\mu}_j(X)) \cdot A / \what{\pi}(X)]}{P_{\rm train} [A / \what{\pi}(X)^2]}$ as in~\eqref{eq:epsilon_xgb}
    \State Compute $\what{g}_{j} = \argmin_{g\in \mathcal G} P_{\rm train}[A( \wt{g}_j - g(X))^2]$ and update  $\what\mu_{j+1} = \what \mu_j - \eta \what{g}_{j}$
\EndFor
\For{$k = 0, 1, \dots, K-1$}
    \State Compute gradient $\wt g_k = \epsilon^\star_{J+k} \cdot 1 / \what{\pi}(X)$ where $\epsilon_{J+k}^\star =  \tfrac{P_{\rm eval} [(Y - \what\mu_{J+k}(X)) \cdot A / \what{\pi}(X)]}{P_{\rm eval}[A / \what{\pi}(X)^2]}$ 
    \State Compute $\what{g}_{k} = \argmin_{g\in \mathcal G} P_{\rm eval}[A( \wt{g}_k - g(X))^2]$ and update  $\what\mu_{J+k+1} = \what \mu_{J+k} - \eta \what{g}_{k}$
\EndFor
\State Return final outcome model $\what\mu^C_{\rm XGB} := \what\mu_{J+K}$
\end{algorithmic}
\label{alg:boosting}
\end{algorithm}


\looseness=-1 \paragraph{Constrained Gradient Boosting}
We want to minimize the squared loss subject to the constraint~\eqref{eq:c-learner} and propose a two-stage procedure. First, we perform gradient boosting where instead of the functional gradient of the loss $g_j$, we use a modified version 
\vspace{-0.3em}
\begin{align}
    \label{eq:epsilon_xgb}
    \wt{g}_{j} := g_{j} + \epsilon^\star_j \cdot \frac{1}{\what{\pi}(X)}~~~\mbox{where}~~~
\epsilon_j^\star = \argmin_\epsilon 
\left\{ P_{\rm train}\left[A\left(Y-\what \mu_j(X)-\epsilon \cdot \frac{A}{\what{\pi}(X)}\right)^2\right]\right\}
    \end{align}
\looseness=-1 given by the targeting objective~\eqref{eqn:targeting} applied to $\what\mu_j$ on the dataset $P_{\rm train}$. 
The modification $\wt g_j$ allows subsequent weak learners to be fit in a direction that reduces the loss \emph{and} makes the plug-in estimator closer to satisfying our constraint on $P_{\rm train}$. To ensure the constraint~\eqref{eq:c-learner-datasets} is satisfied on $P_{\rm eval}$, the second stage fits weak learners  to the gradient of constraint violation below:
\vspace{-0.3em}
\begin{equation*}
    \wt{g}_k := \epsilon^\star_{J+k} \cdot \frac{1}{\what{\pi}(X)}~~\mbox{where}~~
    \epsilon^\star_{J+k} = \argmin_{\epsilon} \left\{ P_{\rm eval}\left[A\left(Y-\what{\mu}_{J+k}(X)-\epsilon \cdot \frac{A}{ \what{\pi}(X)}\right)^2\right]\right\}.
\end{equation*}
We summarize the method above in pseudo-code in Algorithm~\ref{alg:boosting}, which we implement using the XGBoost package with custom objectives~\citep{xgb}. Hyperparameters for the first stage (learning rate $\eta$, and other properties of the weak learners such as max tree depth) are selected to have the lowest MSE loss on $P_{\rm val}$.  Other hyperparameters such as max number of trees $J$ and $K$, may be set on $P_{\rm val}$, or alternatively, by early stopping, with evaluation on different folds within $P_{\rm train}$. 
These hyperparameters are re-used in the second stage. 




\subsection{Neural Networks}
\label{sec:methods_nn}
When outcome models are 
neural networks $\what\mu_\theta(x)$
with weights $\theta$,
we consider the usual MSE loss with a Lagrangian regularizer for the constraint~\eqref{eq:c-learner-datasets}
\vspace{-0.1cm}
\begin{align*}
    L(\theta) = P_{\rm train}\left[A(Y-\what\mu_\theta(X))^2\right] +
\lambda\cdot P_{\rm eval}\left[\frac{A}{\what\pi(X)}(Y-\what\mu_\theta(X))\right]^2.
\vspace{-0.3cm}
\end{align*}
\looseness=-1 We optimize the objective using stochastic gradient methods, where we take mini-batches of the training set to approximate the gradient of the first term and take a full-batch gradient on the evaluation set for the second term. At the end of every training epoch, the constraint on $P_{\rm eval}$ is enforced \emph{exactly} by
adjusting the constant bias term $\theta_{\rm bias}$ in the neural network:
\vspace{-0.1cm}
\begin{align}
\label{eq:const_shift}
    \theta_{\rm bias} \gets \theta_{\rm bias} +
    \left( P_{\rm eval}\left[\frac{A}{\what\pi(X)}\right] \right)^{-1} P_{\rm eval}\left[\frac{A}{\what\pi(X)}(Y-\what\mu_\theta(X))\right].
    \vspace{-0.3cm}
\end{align}
\looseness=-1 We choose to stop training at the epoch that minimizes $P_{\rm val}  \left[ A(Y-\what\mu_\theta(X))^2\right]$, the MSE loss on the validation split. We consider two options for choosing  hyperparameters (e.g. learning rate, $\lambda$).
The first option is to minimize MSE on $P_{\rm val}$. Second,
 since a small bias shift indicates a regularizer that is successful,
 we choose
  among hyperparameters with reasonable MSE loss on $P_{\rm val}$ the one that minimizes the size of the bias shift \eqref{eq:const_shift} in the first epoch. 
While the first method is more standard, the second encourages satisfying the constraint over the course of training, to avoid big jumps in optimization. 
Model selection for nuisance parameters is an active area of research \citep{hahn2019bayesian,rolling2013,setodji2017right}; we explore these in Section~\ref{sec:civilcomments}. 
